\documentclass[a4paper,10pt]{article}

\usepackage[T1]{fontenc}
\usepackage{xcolor}
\usepackage{geometry}
\usepackage{enumitem}
\usepackage{graphicx}
\usepackage{fontawesome5}
\usepackage{mathpazo}
\usepackage{hyperref}

\definecolor{darkblue}{RGB}{0,0,139}
\geometry{left=1.25cm,top=1.2cm,right=1.05cm,bottom=0.8cm}
\pagestyle{empty}
\hypersetup{colorlinks=true,linkcolor=darkblue,filecolor=darkblue,urlcolor=darkblue,
  pdftitle={Flavio D. Gualtieri --- CV},pdfauthor={Flavio D. Gualtieri}}
\urlstyle{same}
\raggedright
\setlength{\tabcolsep}{0pt}

%----------------------------------------------------------------------------------------
% SIDE SECTION LAYOUT
%----------------------------------------------------------------------------------------

\newlength{\sectionlabelwidth}
\newlength{\sectionlabelsep}
\setlength{\sectionlabelwidth}{2.2cm}
\setlength{\sectionlabelsep}{0.35cm}

\newenvironment{sideSection}[1]{
  \vspace{2.2mm}
  \begin{list}{}{
    \setlength{\leftmargin}{\dimexpr\sectionlabelwidth+\sectionlabelsep\relax}
    \setlength{\labelwidth}{\sectionlabelwidth}
    \setlength{\labelsep}{\sectionlabelsep}
    \setlength{\topsep}{0pt}
    \setlength{\partopsep}{0pt}
    \setlength{\parsep}{0pt}
    \setlength{\itemsep}{0pt}
  }
  % zero depth, so a two-line label doesn't push down the first entry
  \item[{\raisebox{0pt}[\height][0pt]{\parbox[t]{\sectionlabelwidth}{\raggedright\bfseries\small\MakeUppercase{#1}}}}]
  \vspace{-0.6mm}
}{
  \end{list}
}

%----------------------------------------------------------------------------------------
% CUSTOM COMMANDS
%----------------------------------------------------------------------------------------

\newcommand{\resumeSubHeadingListEnd}{\end{itemize}}
\newcommand{\FILL}[1]{\textcolor{red}{\textbf{[#1]}}}
\newcommand{\cvline}[1]{\vspace{0.3mm}{\small #1}\vspace{0.6mm}}
\newcommand{\socialicon}[1]{\raisebox{-0.05em}{\resizebox{!}{1em}{#1}}}
\newcommand*{\bigcdot}{\mathbin{\vcenter{\hbox{\scalebox{2}{$\cdot$}}}}}\newcommand{\resumeSubHeadingListStart}{\begin{itemize}[leftmargin=*,labelsep=1mm,label={}]}

\newcommand{\resumeSubheading}[4]{
  \vspace{0.5mm}\item
  \begin{tabular*}{\linewidth}[t]{l@{\extracolsep{\fill}}r}
    \textbf{#1} & \textit{\footnotesize #4} \\
    \textit{\footnotesize #3} & \footnotesize #2 \\
  \end{tabular*}
  \vspace{-1.4mm}
}

\newcommand{\resumeProject}[2]{
  \vspace{0.5mm}\item
  \begin{tabular*}{\linewidth}[t]{l@{\extracolsep{\fill}}r}
    \textbf{#1} & \textit{\footnotesize #2}
  \end{tabular*}\par
}

\newcommand{\cvbullet}[1]{%
  {\small\begin{itemize}[leftmargin=0.9em,labelwidth=0.6em,labelsep=0.3em,align=left,
    label=$\circ$,topsep=0.4mm,partopsep=0pt,parsep=0pt,itemsep=0pt]
  \item #1
  \end{itemize}}}

\begin{document}
\fontfamily{ppl}\selectfont

%----------------------------------------------------------------------------------------
% HEADER
%----------------------------------------------------------------------------------------

\begin{center}
  {\Large\textbf{FLAVIO D. GUALTIERI}}
\end{center}
\vspace{-4mm}

\begin{center}
  \small
  {\textcolor{darkblue}{\socialicon{\faGithub}}}\ \href{https://github.com/flavio-gualtieri}{flavio-gualtieri} $\bigcdot$
  \href{mailto:f.d.gualtieri@qmul.ac.uk}{f.d.gualtieri@qmul.ac.uk} $\bigcdot$
  \href{tel:+447767366389}{+44 7767 366389}
\end{center}
\vspace{-2mm}

%----------------------------------------------------------------------------------------
% EDUCATION
%----------------------------------------------------------------------------------------

\begin{sideSection}{Education}
\resumeSubHeadingListStart

  \resumeSubheading
  {Queen Mary University of London}
  {London, UK}
  {Ph.D. in Mathematics $\bigcdot$ Supervisors: Omer Bobrowski \& Primo\v{z} Skraba $\bigcdot$ Funded by the \href{https://erlangenhub.ox.ac.uk}{Erlangen AI Hub}}
  {Dec 2025 -- Dec 2028}

  \cvline{Researching deep learning models for generation and inference on stochastic point processes.}

  %--------------------------------------------------------------------------

  \resumeSubheading
  {King's College London}
  {London, UK}
  {M.Sc. in Theoretical Physics}
  {Sep 2023 -- Sep 2024}

  \cvline{\emph{Coursework}: quantum field theory, conformal field theory, Lie algebras, stochastic processes. 
  \\ \emph{Dissertation}: \href{https://flavio-gualtieri.github.io/msc-thesis/}{Topological Surface Codes on Exactly Solvable Lattice Models}.}
  
  %--------------------------------------------------------------------------

  \resumeSubheading
  {University of Chicago}
  {Chicago, USA}
  {B.S. in Mathematics \& B.A. in Physics}
  {Sep 2019 -- Jun 2023}

  \cvline{\emph{Coursework}: 
  abstract algebra,
  linear algebra,
  martingales and Brownian motion,
  statistical models and methods,
  statistical mechanics,
  advanced quantum mechanics,
  experimental physics.}

\resumeSubHeadingListEnd
\end{sideSection}

%----------------------------------------------------------------------------------------
% RESEARCH & PROJECTS
%----------------------------------------------------------------------------------------

\begin{sideSection}{Research \& Projects}
\resumeSubHeadingListStart

  %--------------------------------------------------------------------------

  \resumeProject
  {Simulation-Based Inference for Spatial Point Processes}
  {Likelihood-Free Inference, Classification}

  \cvbullet{From one observed point pattern, a classifier picks the generating
  family and a per-family neural estimator infers its parameters, trained on 900k simulations.
  Inputs combine summary functions and topological features.}

  \cvbullet{Designed evaluation method: simulate new patterns using the model's output and score against the input
  with a kernel score on local configurations, giving a score from 0 (random fit) to 1 (true model). The model
  scores 0.90, against 0.74 for classical baselines.}

  %--------------------------------------------------------------------------

  \resumeProject
  {Eigengrokking: Can Hessian Curvature Forecast Grokking?
  \href{https://github.com/flavio-gualtieri/eigengrokking}{\textcolor{darkblue}{\socialicon{\faGithub}}}
  }
  {Transformers, Training Dynamics}

  \cvbullet{Trained 240 transformers on modular arithmetic, tracked Hessian and Gauss--Newton spectra
  through training, and forecast each run's grokking time from the first 500 steps after
  memorization. Built a stochastic Lanczos quadrature pipeline, validated against exact Hessians
  on smaller models.}

  \cvbullet{Found that curvature adds only a modest improvement over simple training
  signals (loss, weight and gradient norms) in explaining why identically configured runs
  grok at different times: $R^2$ increases 0.17$\to$0.28 (ridge regression) and 0.48$\to$0.60 (gradient boosted trees)
  on unseen configurations. The sharpest spectral signal is a collapse in effective rank,
  but it coincides with grokking rather than preceding it. Traced post-grokking accuracy crashes to
  AdamW's default $\beta_2 = 0.999$; applying $\beta_2 = 0.98$ cut them tenfold.}


  %--------------------------------------------------------------------------

  \resumeProject
  {DeepGARCH: Feature-Conditioned Neural Volatility Model
  \href{https://github.com/flavio-gualtieri/deepgarch}{\textcolor{darkblue}{\socialicon{\faGithub}}}}
  {Volatility Forecasting, Time Series}

  \cvbullet{Built a natural gas volatility forecaster in which a neural network (MLP) updates a GARCH(1,1)
  model's parameters each day from recent returns, realized volatility, volume and EIA storage 
  reports, trained end to end. Over 10 training runs it had the lowest forecast error (QLIKE) of five 
  models, ahead of GARCH and EGARCH, though not by a significant margin (Diebold--Mariano $p = 0.11$). 
  The edge depends on tightly capping the model's long-run volatility level, which matters after the 2021 
  gas-price shock but cannot be detected on pre-2021 validation data.}

\resumeSubHeadingListEnd
\end{sideSection}

%----------------------------------------------------------------------------------------
% EXPERIENCE
%----------------------------------------------------------------------------------------

\begin{sideSection}{Experience}
\resumeSubHeadingListStart

  \resumeSubheading
  {PwC Data \& AI}
  {Milan, Italy}
  {Data Scientist (Associate)}
  {Dec 2024 -- Dec 2025}

  \cvbullet{Built an outage-detection pipeline, deployed on the client's AWS, processing 5 GB of telemetry
  per day. An XGBoost time-series model forecasts expected live user counts; a drop below the forecast flags a possible
  outage and triggers DBSCAN clustering on the locations of recently disconnected users. Cluster sizes are
  compared with each area's historical norm, with thresholds tuned by grid search, to separate genuine localized
  outages from background noise. Cut detection delay from 6 h to 30 min.}

  \cvbullet{Engineered new categorical and numerical features for a probabilistic deep-learning model forecasting shipment
  arrival times: a dense neural network (TensorFlow) outputs Weibull shape and scale parameters, trained with a negative
  log-likelihood loss, giving a full distribution over arrival times. Tuned hyperparameters
  with Optuna to 85\% held-out coverage.
  Presented SHAP feature importances to the client for model explainability.}

  %--------------------------------------------------------------------------

  \resumeSubheading
  {Center for RISC, University of Chicago}
  {Chicago, USA}
  {Data Science Intern}
  {Sep 2022 -- May 2023}

  \cvbullet{Built the declustering stage of a pipeline auditing a geospatial-intelligence
  firm's public earthquake predictions against live USGS data. Declustering reduces each
  seismic sequence to one `mainshock', so predictions cannot score a hit on an aftershock by chance.
  Designed a rule-based clustering method:
  events are grouped only if \emph{every} pair lies within 50 miles and 3 days, and the largest-magnitude event in each group is
  the mainshock.}

\resumeSubHeadingListEnd
\end{sideSection}

%----------------------------------------------------------------------------------------
% SKILLS & LANGUAGES
%----------------------------------------------------------------------------------------

\begin{sideSection}{Skills \& Languages}
{\small
\begin{tabular}[t]{@{}p{\linewidth}@{}}
\textbf{Programming:}
Python [\emph{PyTorch, TensorFlow, statsmodels, arch, XGBoost, pandas, NumPy, GUDHI}],
SQL,
R
\\[0.5mm]
\textbf{Quantitative:}
machine learning,
neural networks,
diffusion models,
gradient boosting,
statistical inference,
anomaly detection,
time-series forecasting,
Monte Carlo methods,
stochastic point processes,
topological data analysis
\\[0.5mm]
\textbf{Compute:} Git, AWS [\emph{EC2, S3}], Slurm HPC\\[0.5mm]
\textbf{Languages:} English \& Italian [\emph{native}]\\[0.5mm]
\textbf{Interests:} world capitals
\end{tabular}
}
\end{sideSection}

\end{document}